\section{Assembling prescribed block embeddings}
\label{rep:section}

This section gives an exact assembly criterion for arbitrary common graphs.
The criterion takes the embeddings of the blocks of the two input graphs as
part of its input.  It does not choose those embeddings.  Its purpose is to
separate that remaining choice from the problem of assembling blocks and
placing disconnected common components.  For a common cycle the criterion
reduces to the side condition in the repair theorem of the cactus
manuscript~\cite{cactus}.  For a general common block it contains an additional,
necessary requirement: the chosen face must be incident with every attachment
of the component placed there.

\subsection{The face-list criterion}
\label{rep:criterion}

Throughout this section, $G_1$ and $G_2$ are connected simple planar graphs on
the same vertex set, and their union $U=G_1\cup G_2$ is biconnected and has at
least three vertices.  Their common graph $H=G_1\cap G_2$ is arbitrary.  In
particular, $H$ may have bridges, isolated vertices, and disconnected
components.  A \emph{block} with an edge is either a maximal biconnected
subgraph with at least three vertices or a bridge together with its ends.
Blocks of different colours are treated as different objects, even when their
edge sets coincide.

Prescribe an oriented planar embedding $\mathcal E_B$ of every block $B$ of
each colour.  Write $\rho_B(v)$ for its rotation at $v$.  Embeddings are on an
oriented sphere; no outer face is distinguished, and a prescribed embedding
is not silently allowed to be reflected.  Every common block $K$ is contained
in a unique block $B_i(K)$ of colour $i$.  For a common bridge this is the
unique colour block containing its edge.

For a vertex $v$, let $D_H(v)$ be its set of common darts, and let $D_B(v)$
be its dart set in a colour block $B$.  The first condition
on the prescribed embeddings is
\begin{equation}
\tag{O}\label{rep:eq:order}
\begin{gathered}
\text{there is a cyclic order $\omega_v$ of $D_H(v)$ that restricts to}
\\[-1mm]
\rho_B(v)|_{D_H(v)\cap D_B(v)}\text{ for every colour block $B$ at $v$,}
\\[-1mm]
\text{and makes the partition by blocks of each colour noncrossing.}
\end{gathered}
\end{equation}
Empty parts are ignored.  Here two disjoint parts interleave if four distinct
elements, two from each part, occur alternately in the cyclic order.

If~\eqref{rep:eq:order} holds, the embeddings induced on a common block $K$
from $\mathcal E_{B_1(K)}$ and $\mathcal E_{B_2(K)}$ have the same rotation at
every vertex of $K$.  They therefore define the same oriented embedding of
the connected graph $K$.  Denote it by $\kappa_K$ and its set of faces by
$\mathcal F(K)$.  Faces may be identified by their oriented facial walks.
The two faces of a cycle remain distinct under this convention, although
their unoriented boundary edge sets are equal.  A bridge block has one face.
The faces $F_i(K,Y)$ and the entries of the lists $L(K,Z)$ below are faces of $\kappa_K$, so
they are defined only when~\eqref{rep:eq:order} holds.  Accordingly, in
Theorem~\ref{rep:thm:repair} the condition $L(K,Z)\ne\varnothing$ is
evaluated under~\eqref{rep:eq:order}, and the algorithm of
Section~\ref{rep:algorithm} forms the lists only after
checking~\eqref{rep:eq:order}.

For $Z\in\cc(U-V(K))$, put
\[
 A_K(Z)=N_U(Z)\cap V(K).
\]
A \emph{colour-$i$ lobe of $K$} is a component
$Y\in\cc(B_i(K)-V(K))$.  Its prescribed embedding places $Y$ in a unique
face of $\kappa_K$, denoted by $F_i(K,Y)$.  Every such $Y$ is contained in a
unique component $Z$ of $U-V(K)$.  Define
\begin{equation}
\label{rep:eq:face-list}
\begin{split}
 L(K,Z)={}&\{F\in\mathcal F(K):A_K(Z)\subseteq V(\partial F)\}
 \\
 &{}\cap
 \bigcap_{\substack{i\in\{1,2\},\;Y\in\cc(B_i(K)-V(K))\\Y\subseteq Z}}
 \{F_i(K,Y)\}.
\end{split}
\end{equation}
If there is no lobe in the displayed intersection, only the attachment
condition in the first line is imposed.

\begin{theorem}[Face-list repair]
\label{rep:thm:repair}
There is a \SEFE{} whose colour embeddings restrict to every prescribed
$\mathcal E_B$ if and only if~\eqref{rep:eq:order} holds and
$L(K,Z)\ne\varnothing$ for every common block $K$ and every
$Z\in\cc(U-V(K))$.

Moreover, after independently choosing
\[
 F(K,Z)\in L(K,Z)
\]
for every such pair, one can construct a \SEFE{} that preserves all
$\mathcal E_B$ and places $Z$ in $F(K,Z)$ relative to $K$.  The conditions can
be checked, and compatible combinatorial embeddings constructed, in
deterministic $O(N^3)$ time.
\end{theorem}

For a common cycle, both faces contain every attachment vertex.  Thus its
attachment condition is automatic, and the face-list test asks exactly that
the prescribed lobe sides agree on each union component.  For a common
bridge, its unique face contains both ends, and its whole face-list condition
is automatic.  The theorem therefore extends the structural repair step
without a bridgelessness assumption on $H$.

\subsection{Two elementary embedding facts}
\label{rep:elementary}

We first record how to check~\eqref{rep:eq:order} and how block faces distinguish
faces of a connected common component.

\begin{lemma}[Two noncrossing partitions]
\label{rep:lem:two-partitions}
Let a finite set $D$ have a red partition and a blue partition, with a cyclic
order $\sigma_A$ prescribed on every part $A$.  Form the bipartite multigraph
$M$ whose vertices are the parts and whose edge labelled $d$ joins the two
parts containing $d$.  Give a red vertex $A$ rotation $\sigma_A$ and a blue
vertex $A$ rotation $\rev\sigma_A$.

This rotation system is planar if and only if there is a cyclic order of $D$
that restricts to every $\sigma_A$ and makes both partitions noncrossing.
Such an order can be constructed in $O(|D|)$ time.
\end{lemma}

\begin{proof}
Suppose the cyclic order exists.  Place its elements on a circle.  Draw the
red parts as disjoint stars in the interior and the blue parts as disjoint
stars in the exterior.  Noninterleaving parts admit these disjoint stars:
one may use their convex hulls on each side, with arbitrarily small
perturbations at boundary points.  Joining the two half-edges at each labelled
point gives the required planar rotation system of $M$; the exterior stars
have the reversed cyclic order.

Conversely, consider one connected component of the prescribed plane map
$M$, and subdivide every edge $d$ at a point $m_d$.  At each corner of a blue
vertex, draw an arc between the red ends of its two incident edges, running
along those edges and through the corner.  The arcs can be chosen mutually
disjoint except at red endpoints: use disjoint narrow ribbons beside the
edges and disjoint small sectors inside each corner.  Around a red endpoint,
an arc using an edge leaves on its specified side, so these choices are
consistent there as well.  The auxiliary multigraph formed by these arcs on
the red vertices is connected, because the arcs around each blue vertex
connect all its red neighbours and $M$ is connected.  Choose a spanning tree
of these arcs.

Together with the red halves of all original edges, the chosen arcs form an
embedded tree $T$ whose leaves are the points $m_d$.  Take a sufficiently thin
closed neighbourhood of $T$.  It is a disk, and its boundary crosses each
original edge exactly once near $m_d$ and no other original edge.  All red
stars lie inside this disk and all blue stars lie outside.  The cyclic order
of the crossing labels therefore restricts to the red rotations and to the
reverses of the blue rotations.  Stars on the same side cannot interleave:
two interleaving pairs of boundary attachments would force their connecting
arcs to cross.  This is the desired order.  For a disconnected $M$, concatenate
the orders obtained for its components.

The construction uses one auxiliary arc per blue corner, a spanning-tree
search, and one traversal of the resulting tree.  All have linear size and
take $O(|D|)$ time.  Planarity of the prescribed rotation system is checked
by tracing its face walks and applying Euler's formula separately to each
connected component.
\end{proof}

Apply this lemma at $v$ with the parts given by its colour blocks and the
elements given by its common darts.  The total size of all these maps is
$O(N)$.

\begin{lemma}[Block faces distinguish faces]
\label{rep:lem:face-signatures}
Let $X$ be a connected plane graph.  A face of $X$ is uniquely determined by
the faces containing it in the embeddings of the blocks of $X$ that have
edges.  Bridge coordinates are harmless: each has only one value.  An
isolated vertex has just one face.
\end{lemma}

\begin{proof}
It suffices to show that two distinct faces $f,f'$ of $X$ are separated by a
simple cycle of $X$.  Choose a curve $\gamma$ from a point of $f$ to a point
of $f'$, transverse to edges and avoiding vertices.  In the binary edge
vector of the facial boundary of $f$, edges occurring twice cancel.  This
vector has even degree at every vertex and hence decomposes into simple
cycles.  Its intersection parity with $\gamma$ is one: each crossing of an
edge incident with $f$ on exactly one side changes whether the curve is in
$f$, and the curve starts in $f$ and ends outside it.  Therefore at least one
cycle in the decomposition has odd intersection parity with $\gamma$.
The Jordan curve theorem says that this cycle separates $f$ from $f'$.

Every simple cycle belongs to a single block $K$ of $X$.  Two regions in the
same face of $K$ lie on the same side of every cycle of $K$, since that face
is connected and avoids the cycle.  Hence the separating cycle implies that
$f$ and $f'$ have different $K$-face coordinates.
\end{proof}

We will also use the following block-assembly fact.  At a cutvertex of a
planar graph, distinct block dart sets do not interleave.  Conversely,
prescribed planar block rotations can be assembled using any noncrossing
cyclic merge at each cutvertex.  To verify the converse, first merge pieces
sharing just one vertex.  A noncrossing partition has an innermost part that
is an interval in the restriction to its parts.  To see this, if no part
were an interval, choose a nonempty gap between consecutive elements of
one part with minimum cardinality.  Every other part meeting that gap
lies wholly in it, or the two parts would interleave.  Such a part, being
noninterval, has a nonempty gap of its own inside the chosen gap, a
contradiction.  Insert the interval piece in a
corner of the others; this glues two disk faces along the chosen corner and
preserves genus zero.  Remove the interval and induct.  Finally induct over
the block-cut tree.  For necessity, two alternating pairs of darts from
different blocks would yield two crossing connections after deleting the
cutvertex, contradicting planarity.  These observations also apply when one
of the pieces is a bridge.

\subsection{The local region tree}
\label{rep:signature-tree}

Assume~\eqref{rep:eq:order} and choose $F(K,Z)\in L(K,Z)$ for every pair.
For a connected subgraph $Y$ disjoint from $K$, write $F(K,Y)$ for the
selected face of the component of $U-V(K)$ containing $Y$.

Fix a vertex $v$ incident with a common edge, and let $\mathcal K_v$ be
the common blocks containing $v$.  Distinct members meet only at $v$,
and $K-v$ is connected for every $K\in\mathcal K_v$.  For each such $K$,
let $\mathcal F_v(K)$ consist of the faces of $\kappa_K$ incident with $v$.
Each of these faces has a unique corner at $v$; for a bridge its unique
face has its unique corner.  In particular,
$|\mathcal F_v(K)|=\deg_K(v)$.

We construct a tree from these corners.  Define an adjacency graph $A_v$
on $\mathcal K_v$: two blocks $K,J$ are adjacent when a path in $U-v$
joins $K-v$ to $J-v$ and has no internal vertex in any other member of
$\mathcal K_v$.  Since $U-v$ is connected, $A_v$ is connected.  If $KJ$
is an edge of $A_v$ and $M$ is a third common block at $v$, the two
block interiors and the connecting path avoid $M$.  Hence
\begin{equation}\label{rep:eq:adjacent-face}
 F(M,K-v)=F(M,J-v).
\end{equation}

For each $K$ construct a star with centre $K$ and one \emph{port}~$g$
for every $g\in\mathcal F_v(K)$; it represents the unique corner of $g$
at~$v$.  Choose an arbitrary spanning tree
$T^0_v$ of $A_v$.  For each edge $KJ$ of $T^0_v$, identify the port
$F(K,J-v)$ of $K$ with the port $F(J,K-v)$ of $J$.  These ports exist:
the indicated interiors are adjacent to $v$, so the attachment condition
in~\eqref{rep:eq:face-list} makes the selected faces incident with $v$.
Let $T_v$ denote the resulting bipartite block--region graph.

\begin{lemma}[Region tree and its face coordinates]
\label{rep:lem:signature-tree}
The graph $T_v$ is a tree.  It does not depend, up to its named
incidences, on the choice of $T^0_v$.  Every region-node $t$ has a
well-defined vector $q(t)$ with one $J$-face coordinate for each
$J\in\mathcal K_v$.  The incidence at the $g$-port of $K$ has vector
\begin{equation}\label{rep:eq:signature}
 q(K,g)_J=\begin{cases}
          g,&J=K,\\
          F(J,K-v),&J\ne K.
         \end{cases}
\end{equation}
Distinct region-nodes have distinct vectors.  After deleting a block-node
$K$, its component through the $g$-port consists precisely of nodes with
$K$-coordinate $g$.
\end{lemma}

\begin{proof}
Root $T^0_v$ and add its corner-stars in parent-before-child order.  A
new star meets the existing construction at exactly one port, so the
result is a tree.

For an edge $KJ$ of $T^0_v$, the two ports identified by the
construction have the same vector in~\eqref{rep:eq:signature}: their
$K$- and $J$-coordinates agree by the definition of those ports, and
all third coordinates agree by~\eqref{rep:eq:adjacent-face}.  Thus the
vectors descend to region-nodes of $T_v$.

Every incidence not meeting $K$ preserves the $K$-coordinate.  The
neighbours of $K$ have distinct $K$-coordinates, because the ports of
$K$ are its distinct corners.  Since $T_v$ is a tree, deleting $K$
leaves one component per port, and the indicated coordinate is constant
on each component.  If two distinct region-nodes had equal coordinate
vectors, their unique connecting path would traverse some block-node
$K$ through two different ports, forcing different $K$-coordinates.

Finally consider the graph formed by all corner-stars with ports
identified whenever their vectors agree.  Its distinct region-nodes
remain separated by the preceding argument, and its named incidence
edges are those of $T_v$.  Hence it is the same tree and is independent
of the chosen spanning tree $T^0_v$.
\end{proof}

The coordinates are only a way to \emph{identify} the regions: the tree
itself is constructed by a spanning-tree gluing, not by solving a system
of face-signature equations.

\begin{lemma}[Locating a block without common darts]
\label{rep:lem:unmarked-block-signature}
Let $Y$ be connected, disjoint from every $K-v$ with
$K\in\mathcal K_v$, and contained in $U-v$.  If $Y$ has an edge to $v$,
then the vector
\[
 q(Y)_K=F(K,Y)\qquad(K\in\mathcal K_v)
\]
labels a region-node of $T_v$.
\end{lemma}
\begin{proof}
Follow a path in $U-v$ from $Y$ to its first meeting with a local
common block, say $K-v$.  This path avoids every other $J-v$, so
$F(J,Y)=F(J,K-v)$ for $J\ne K$.  The edge from $Y$ to $v$ makes
$F(K,Y)$ incident with $v$.  Consequently $q(Y)$ equals the vector
$q(K,F(K,Y))$ of an existing region-node.
\end{proof}

We use the lemma when $Y=B-v$ for a colour-$i$ block $B$ at $v$ with no
common dart there.  Its hypotheses hold.  The graph $B-v$ is connected
and has an edge to $v$, since $B$ is a block with an edge at $v$.  For
$K\in\mathcal K_v$ the block $B_i(K)$ contains the common darts of $K$ at
$v$, so $B\ne B_i(K)$.  Two distinct colour-$i$ blocks share at most one
vertex, here $v$, so $B-v$ is disjoint from $B_i(K)-v\supseteq K-v$.  The
adjacency-to-$v$ hypothesis cannot be dropped for an arbitrary connected
subgraph of $U-v$.

\subsection{Plane bouquets and the contour representation}
\label{rep:bouquets}

A \emph{block bouquet at $v$} is a union $W$ of blocks
$K_1,\ldots,K_t$ that pairwise meet exactly at $v$.  Each block is either
biconnected or a single edge.  Give $W$ a planar embedding, and consider only
its faces incident with $v$.

For a face $g$ of a biconnected $K_j$ incident with $v$, the facial boundary
minus $v$ is a connected path disjoint from every other block.  A sufficiently
thin strip along the $g$-side of that path lies in a single face $f$ of $W$.
Both ends of the strip approach the corner $g$ at $v$, so $f$ is incident
with $v$.  This defines the global face adjacent to the port $g$.  For a
single-edge block, the edge is dangling in the bouquet, and its two sides
belong to one face; its unique port is assigned to that face.

Let $S(W,v)$ be the bipartite graph whose nodes are the blocks $K_j$ and the
faces of $W$ incident with $v$, with one incidence from $K_j$ to the face just
defined for each port $g\in\mathcal F_v(K_j)$.

\begin{lemma}[Bouquet contour]
\label{rep:lem:bouquet}
The graph $S(W,v)$ is a tree.  Its region-nodes have distinct vectors of faces
of the individual blocks.  Give each block-node $K$ the cyclic order of its
corners at $v$.  There is exactly one choice of rotations at the region-nodes
for which the contour of this plane tree recovers the rotation of $W$ at $v$
as follows: label the outgoing traversal from $K$ through a port $g$ by the
dart immediately preceding that corner in the rotation of $K$.

Furthermore, choose any distinguished dart $d_K$ of every block at $v$.
At a region-node $t$, its rotation is the restriction of the resulting dart
order to the representatives $d_K$ for the block-nodes adjacent to $t$.
\end{lemma}

\begin{proof}
Distinct faces of $W$ have distinct block-face vectors by
Lemma~\ref{rep:lem:face-signatures}.  All ports of a fixed block $K$ have the
same coordinate for every other block $J$: it is the face of $J$ containing
$K-v$.  Their $K$-coordinates are their respective ports.  Hence a simple
cycle in $S(W,v)$ would change one block coordinate exactly once, by the
same argument as in Lemma~\ref{rep:lem:signature-tree}.  Thus $S(W,v)$ is a
forest.

Take a small positively oriented circle around $v$.  On crossing a dart of
$K$, move in $S(W,v)$ from the region before that dart to $K$ and then to the
region after it.  These are exactly the two ports bordering that dart in
the rotation of $K$.  The strip construction above ensures that the region
on each side is the one assigned to the corresponding port, even if other
blocks have darts between consecutive $K$-darts.  As the circle makes one
turn, it gives a single closed walk using every incidence once in each
direction.  Every region-node is visited, because a face incident with $v$
has a sector bordered by a dart at $v$.  Thus $S(W,v)$ is connected and is
a tree.

The circle walk prescribes a successor for each incoming incidence at every
region-node.  This successor permutation has a single cycle.  Otherwise,
splitting that node according to its several successor cycles would
disconnect the tree, while the circle walk would remain one closed walk
using all its directed incidences, which is impossible.  Hence these
successors are rotations.  At a block-node the successor is the prescribed
cyclic order of its corners.  The walk is therefore precisely the contour
of this plane tree, and the indicated labels are the original darts in
their original order.  The walk determines every successor, proving
uniqueness.

Finally, fix a region-node $t$ and a neighbouring block-node $K$.  The contour
excursion beginning with $t\to K$ and ending with $K\to t$ contains every
outgoing traversal from $K$, including that final traversal.  Thus it
contains all labels of $K$, including the chosen $d_K$, and no labels of
another block-node adjacent to $t$.  The order of the representatives is
therefore exactly the order of the excursions at $t$.
\end{proof}

For a colour block $B$ at $v$, put
\[
 J_v(B)=\{K\in\mathcal K_v:K\subseteq B\}.
\]
These are exactly the common blocks contributing darts of $B$ at $v$.
If $K\notin J_v(B)$, then $B-v$ is connected and disjoint from $K$.
When $J_v(B)\ne\varnothing$, restrict $\mathcal E_B$ to the bouquet
\[
 W_B=\bigcup_{K\in J_v(B)}K
\]
and let $S_B=S(W_B,v)$ be its plane incidence tree from
Lemma~\ref{rep:lem:bouquet}.

\begin{lemma}[Colour blocks occupy subtrees]
\label{rep:lem:block-subtree}
The tree $S_B$ is naturally isomorphic to the subtree $T_B$ of $T_v$ consisting
of all incidences at the block-nodes in $J_v(B)$, together with their ends.
Every region-node of $T_B$ has $K$-coordinate $F(K,B-v)$ whenever
$K\notin J_v(B)$.
\end{lemma}

\begin{proof}
Map each block-node to the identically named block-node of $T_v$.  For a
region $f$ of $S_B$, define a full signature by taking, for
$K\in J_v(B)$, the face of $K$ containing $f$, and, for
$K\notin J_v(B)$, the value $F(K,B-v)$.

Consider the port $g$ of a block $J\in J_v(B)$.  For
$K\in J_v(B)\setminus\{J\}$, points in the boundary strip of that port
lie in the face of $K$ containing $J-v$.  The connected graph $J-v$ is
contained in $B-V(K)$, hence in one lobe of $K$ in this colour.  The
singleton conditions defining $L(K,Z)$ therefore make that face equal to
$F(K,J-v)$.  For $K\notin J_v(B)$, the set $B-v$ is connected and disjoint
from $K$, so $F(K,B-v)=F(K,J-v)$.  The image of the port is consequently
exactly $q(J,g)$.

Distinct region-nodes of $S_B$ have distinct face coordinates among
$J_v(B)$, by Lemma~\ref{rep:lem:bouquet}.  The map is therefore injective
on region-nodes, as well as on block-nodes, and maps incidences to
incidences.  Its image contains exactly all incidences at the nodes in
$J_v(B)$ and their ends.  It is a connected subgraph of the tree $T_v$,
so it is the stated subtree.  The constant-coordinate assertion follows
from the definition of the map.
\end{proof}

\subsection{Constructing the repaired order at a vertex}
\label{rep:local-construction}

Choose an order $\omega_v$ witnessing~\eqref{rep:eq:order}, and choose one
distinguished common dart $d_K$ of each $K\in\mathcal K_v$.  Give the
block-node $K$ of $T_v$ the corner order of $\kappa_K$ at $v$.  At a
region-node $t$, give its incident block-nodes the restriction of
$\omega_v$ to their representatives $d_K$.  Let $\omega'_v$ be the cyclic
order obtained from the contour of this plane tree by labelling an outgoing
traversal $K\to q(K,g)$ with the common dart immediately preceding $g$ in
$\kappa_K$.  This is a cyclic order of exactly $D_H(v)$.  The contour
traverses each incidence $K\to q(K,g)$ exactly once, and the corners of $K$
at $v$ correspond bijectively to the darts of $K$ at $v$, each corner to the
dart preceding it.  Every common dart at $v$ belongs to exactly one block
in $\mathcal K_v$.

\begin{lemma}[Properties of the repaired common order]
\label{rep:lem:common-order}
The order $\omega'_v$ has the following properties.
\begin{enumerate}
\item For every colour block $B$ at $v$, its restriction to the common darts
of $B$ equals the restriction of $\rho_B(v)$.
\item For each colour, the common dart sets of its distinct blocks at $v$
are noncrossing in $\omega'_v$.
\item If $K\in\mathcal K_v\setminus J_v(B)$, all common darts of $B$
lie in the corner of $K$ corresponding to $F(K,B-v)$.
\end{enumerate}
\end{lemma}

\begin{proof}
For the first claim, use the subtree $T_B$ of
Lemma~\ref{rep:lem:block-subtree}.  Its block-node rotations agree with
those of $S_B$.  At a region-node, the order of its neighbouring block-node
representatives obtained from $\omega_v$ equals their order obtained from
$\rho_B(v)$, since~\eqref{rep:eq:order} gives equality on all the common
darts of $B$.  By the representative assertion in
Lemma~\ref{rep:lem:bouquet}, this is precisely the region-node rotation of
$S_B$.  Thus $T_B$, with the inherited rotations, is the plane tree $S_B$.
Restricting a tree contour to a connected subtree deletes the excursions
outside that subtree and leaves its contour.  Also, every incidence at a
block-node of $T_B$ belongs to $T_B$, so no dart label of that block-node is
lost.  Lemma~\ref{rep:lem:bouquet} now gives the first claim.

For the second claim, let $B,B'$ be distinct blocks of one colour, both
having common darts at $v$.  Their subtrees $T_B,T_{B'}$ have disjoint
incidence edges and disjoint block-nodes.  As subtrees of a tree they meet
in at most one node, necessarily a region-node.  If they are disjoint,
consider the first edge on the path from $T_B$ towards $T_{B'}$.  The
contour excursion beyond this edge contains all labels of $B'$ and no
label of $B$, so the two sets do not interleave.  If they meet at a
region-node $t$, their labels occur in the excursions through two disjoint
sets of neighbours of $t$.  These neighbour sets do not interleave in the
rotation at $t$: their representatives belong to the common dart sets of
$B,B'$, which do not interleave in $\omega_v$.  Replacing each neighbour
by its contour excursion cannot introduce interleaving.  More explicitly,
an interleaving of labels in disjoint excursion intervals would yield an
interleaving of their four excursion indices, contradicting the preceding
statement.

For the third claim, Lemma~\ref{rep:lem:block-subtree} puts $T_B$ in the
component of $T_v-K$ whose $K$-coordinate is $F(K,B-v)$.  By the contour
rule, all labels in that component occur between the two consecutive
$K$-darts bordering that corner.  If $K$ is a bridge, this is its unique
corner and the statement has the same meaning.
\end{proof}

The common order must now be extended to all darts without changing any
prescribed block rotation.  We state the elementary insertion argument
explicitly.

\begin{lemma}[Insertion of private runs]
\label{rep:lem:private-runs}
Let disjoint finite sets $D_B$ carry cyclic orders $\theta_B$, and let some
elements be marked.  Suppose a cyclic order $\omega$ of the marked elements
restricts to each $\theta_B$ on its marked elements and makes the nonempty
marked parts noncrossing.  Replace each marked element $d\in D_B$ by $d$
followed by the run of unmarked elements following $d$ in $\theta_B$ before
the next marked element of $D_B$.  Insert each set with no marked element,
cut arbitrarily into a linear order, after one of these whole runs; several
sets at the same place are concatenated.

The result restricts to every $\theta_B$, and the full sets $D_B$ are
noncrossing.  Every inserted element remains in the gap of $\omega$ after
the marked head of the run to which it was assigned.
\end{lemma}

\begin{proof}
The runs of a marked set occur in the order of its marked heads, giving the
correct restriction.  A set with no marked element is inserted as one
interval, so it cannot interleave with another set.  If two marked sets
interleaved after their runs were expanded, the two witnessing elements
of either set would have to lie in different runs: a run is an interval
containing no element of the other set.  Replacing those runs by their
heads would give four alternating marked heads in $\omega$, a
contradiction.  The assertion about gaps follows directly from the order
of the insertions.
\end{proof}

Apply this lemma separately to each colour, with the common darts marked
and with $\omega'_v$ as their order.  For a colour block $B$ at $v$ having
no common dart there, Lemma~\ref{rep:lem:unmarked-block-signature} gives the
region-node $q(B-v)$.  Choose any incidence
$K\to q(B-v)$, and let $d$ be its contour label.  Insert the rotation
$\rho_B(v)$, cut anywhere, after the entire run headed by $d$.

\begin{lemma}[Complete local rotations]
\label{rep:lem:local-rotations}
This construction gives rotations $\rho_1(v),\rho_2(v)$ with the following
properties.
\begin{enumerate}
\item Each $\rho_i(v)$ restricts to $\rho_B(v)$ for every colour-$i$ block
$B$ at $v$, and distinct such block dart sets do not interleave.
\item Both rotations restrict to $\omega'_v$ on their common darts.
\item For every common block $K$ through $v$ and every colour block $B$
at $v$ not containing $K$, all darts of $B$ enter the face
$F(K,B-v)$ of $K$.
\end{enumerate}
\end{lemma}

\begin{proof}
The first two claims follow from Lemmas~\ref{rep:lem:common-order} and
\ref{rep:lem:private-runs}.  For the third, first suppose $B$ has a common
dart at $v$.  Its common darts lie in the required corner of $K$ by
Lemma~\ref{rep:lem:common-order}.  Every private dart of $B$ is inserted
after one of those common darts and before the next common dart of the
whole order.  No $K$-dart is crossed, so its $K$-corner remains the same.

Now suppose $B$ has no common dart at $v$.  It is inserted after a dart
whose contour traversal ends at $q(B-v)$.  Immediately after that
traversal, the contour is at the region-node with $K$-coordinate
$F(K,B-v)$.  If the labelled dart itself belongs to $K$, its succeeding
corner is exactly that coordinate.  Otherwise its block-node and this
region-node lie in the corresponding component of $T_v-K$, and the same
conclusion follows from Lemma~\ref{rep:lem:signature-tree}.  The run
following the dart contains no common dart and therefore does not change
the corner.  Inserting $B$ after the run puts all of its darts into the
claimed face.
\end{proof}

At a vertex with no common edge, concatenate the rotations of its blocks
in each colour, with arbitrary cuts.  This satisfies the block-assembly
conditions and imposes no common-dart requirement.  Such a vertex lies in
no common block, so the third claim of Lemma~\ref{rep:lem:local-rotations}
is not needed there; the global argument below uses that claim only at
vertices of common blocks.

\subsection{Global reconstruction and relative positions}
\label{rep:global}

\begin{proof}[Proof of Theorem~\ref{rep:thm:repair}, apart from the time bound]
Necessity follows from a \SEFE{}.  Its common rotation at each vertex gives
\eqref{rep:eq:order}.  For a common block $K$ and
$Z\in\cc(U-V(K))$, the drawings of the vertices and edges of $U[Z]$ form a
connected set disjoint from the common drawing of $K$.  Private edges of
different colours may cross one another, but this does not affect that
connectedness.  Neither colour can cross an edge of $K$, since every such
edge belongs to both colours.  Hence all of $Z$ lies in one face $F$ of $K$.
Every attachment edge approaches $K$ from $F$, so
$A_K(Z)\subseteq V(\partial F)$.  Every colour lobe contained in $Z$ lies
in that same face.  Thus $F\in L(K,Z)$.

For sufficiency, make the independent face selections in the theorem and
construct the local rotations above.  By
Lemma~\ref{rep:lem:local-rotations} and the block-assembly fact, they give
planar rotation systems of both connected input graphs.  They preserve
every $\mathcal E_B$, and their restrictions to common darts agree.

We prove that every vertex $x\notin V(K)$ lies in the chosen face
$F(K,x)$ of every common block $K$ in both colour embeddings.  Fix a
colour $i$ and set $B=B_i(K)$.  If $x\in V(B)$, the lobe containing $x$
has that face by the defining singleton condition in
\eqref{rep:eq:face-list}; the reconstruction preserves $\mathcal E_B$.

Otherwise the path in the block-cut tree of $G_i$ from $B$ towards $x$
leaves $B$ at a cutvertex $w$.  The vertex $x$ lies in a component $D$ of
$G_i-w$ disjoint from $V(B)$, and $D$ contains $B'-w$ for the first block
$B'\ne B$ of that branch at $w$.  If $w\notin V(K)$, the connected set
$D\cup\{w\}$ avoids $K$.  It lies in the face of $K$ containing $w$,
which is $F(K,w)=F(K,x)$ by the preceding case and the fact that both
vertices belong to one component of $U-V(K)$.

If $w\in V(K)$, Lemma~\ref{rep:lem:local-rotations} puts every dart of
$B'$ at $w$ into the corner $F(K,B'-w)$.  The connected graph $D$ avoids
$K$, so all of $D$ lies in that face.  Since $D$ contains $x$ and $B'-w$,
this face is $F(K,x)$.  This proves the asserted face position for all
vertices.

It remains to check the embedding of the disconnected common graph,
rather than merely its rotations.  Let $X,X'$ be distinct connected
components of $H$, and choose a vertex $x\in V(X')$.  In each colour,
$X'$ is contained in the face of $X$ containing $x$.  For every block $K$
of $X$, that face has the same $K$-coordinate in the two embeddings, namely
$F(K,x)$.  The rotations of $X$ already agree.  By
Lemma~\ref{rep:lem:face-signatures}, the containing face of $X$ is therefore
the same.  If $X$ is an isolated vertex, its unique face makes the assertion
immediate.  Thus all relative positions of common components agree,
including isolated common vertices.

Rotations and these relative positions specify the oriented embedding of
a finite planar graph.  We recall the standard facts behind this
statement.  (T1) Janiszewski's theorem: if compact sets
$A,A'\subset S^2$ have connected or empty intersection and neither
separates two points $x,y$, then $A\cup A'$ does not separate
them~\cite{newman}.  (T2) The Jordan--Schoenflies theorem: a simple closed
curve bounds a closed disk on each side, and a homeomorphism between two
simple closed curves extends to a homeomorphism of the
sphere~\cite[Chapter~2]{mohar}.  By (T2), an orientation-preserving
homeomorphism between the boundaries of two spheres with the same number
of holes that respects a given bijection of their boundary curves extends
to an orientation-preserving homeomorphism between them: cut both along
disjoint arcs joining the boundary curves in the same pattern, which
reduces the extension to closed disks.

For a component $X$ of $H$ and a component $X''\ne X$, let $f_{X''}(X)$
be the face of $X''$ containing $X$, and put $A_{X''}=S^2\setminus
f_{X''}(X)$.  Each $A_{X''}$ is compact and connected, and its complement
is connected.  For a face $f$ of $X$, put $A_X=S^2\setminus f$.  These
sets form a laminar family: if $X_2$ lies in a face $g$ of $X_1$ with
$g\ne f_{X_1}(X)$, then the connected set $S^2\setminus g$ avoids $X_2$
and contains $X$, so every face of $X_2$ other than $f_{X_2}(X)$ lies in
$g$, and $A_{X_2}\subseteq A_{X_1}$; if neither component separates the
other from $X$, the same argument shows $A_{X_1}\cap A_{X_2}=\varnothing$.
The case of $A_X$ is identical.  The union of the family is therefore a
disjoint union of its maximal members, and (T1) shows that
\[
 R(X,f)=f\cap\bigcap_{X''\ne X}f_{X''}(X)
\]
is connected.  It avoids $H$ and contains the points of $f$ near the
facial walk of $f$, so it is the face of $H$ along that walk.  It follows
that $R(X,f)=R(X',f')$ if and only if $X'\subseteq f$, $X\subseteq f'$,
and no third component separates $X$ from $X'$.  Thus the faces of $H$,
and the facial walks of components bordering each of them, are determined
by the rotations and the pairwise relative positions, which agree in the
two colours.

Now take thin ribbon neighbourhoods of the components in a polygonal model
of each drawing.  Equal rotations identify each component's neighbourhood,
orientation-preservingly, with boundary curves corresponding to faces.  By
the preceding paragraph each complementary region is a sphere with holes
bounded by the same boundary curves in both drawings, and (T2) extends the
boundary identification over it.  Gluing gives an orientation-preserving
homeomorphism of the sphere identifying the two drawings of $H$.  Apply
this homeomorphism to one of the colour drawings.  The resulting drawings
form a \SEFE{} and preserve the prescribed oriented colour-block
embeddings.
\end{proof}

\subsection{Algorithm and complexity}
\label{rep:algorithm}

The proof is constructive.  The following steps specify the resulting
algorithm for the prescribed-block problem.
\begin{enumerate}
\item Compute the block-cut trees of both colours and of $H$.  Check each
prescribed rotation system by face tracing and Euler's formula.
\item At every vertex form the bipartite map of
Lemma~\ref{rep:lem:two-partitions}.  Reject if its prescribed rotation system
is nonplanar; otherwise extract $\omega_v$.
\item Restrict the block embeddings to every common block $K$, trace its
faces, and compute the components of $U-V(K)$, their attachment sets, and
the colour lobes.  Form the lists~\eqref{rep:eq:face-list}.  Reject if any
list is empty; otherwise select any one face from each list.
\item Construct the signature trees $T_v$, their rotations, and their
labelled contours.  Insert private runs and the blocks with no common
dart at $v$ as above.  Output the resulting rotations of $G_1,G_2$.
\end{enumerate}

All face pins are computable from rotations.  For example, to locate a
lobe $Y$ relative to $K$, choose an attachment edge $ay$, with $a\in K$
and $y\in Y$.  In $\rho_{B_i(K)}(a)$, find the preceding $K$-dart before
$ay$.  The corresponding corner of the restricted $K$-rotation identifies
the face containing $Y$.  For a bridge block there is only one face.
Planarity and connectedness of $Y$ guarantee that every attachment gives
the same answer.

For completeness, the number of components and lobes has a useful linear
bound.  Root each component of the block-cut forest of $H$.  For a common
block $K$, let $r_K$ be its parent cutvertex, or choose any vertex of $K$
if it is a root block.  The sets
\[
 V(K)\setminus\{r_K\}
\]
are pairwise disjoint.  Every $Z\in\cc(U-V(K))$ has at least two
attachment vertices, since zero attachments disconnects $U$ and one
attachment is a cutvertex of $U$.  Charge $Z$ to an attachment edge whose
$K$-endpoint is not $r_K$.  An endpoint determines at most one common
block under this rule, and the other endpoint determines the component
$Z$.  Each edge of $U$ is therefore charged at most once per endpoint,
giving
\[
 \sum_K |\cc(U-V(K))|\le 2|E(U)|.
\]
The same argument inside each biconnected colour block counts its lobes
for the common blocks contained in it.  A colour block that is a bridge
has no lobe of its common edge.  Summing over colour blocks gives $O(N)$
lobes in total.

There are $O(N)$ common blocks, so recomputing components and relevant
face information once per block takes $O(N^2)$ time.  Each face list has
at most $O(N)$ entries, and the straightforward attachment-incidence
checks and singleton intersections take at most $O(N^3)$ time overall.
No succinct or implicit face representation is needed for this bound.

At $v$, put $h_v=\deg_H(v)$ and $k_v=|\mathcal K_v|$.  The candidate
signatures $q(K,g)$ have $h_v$ incidences and $k_v$ coordinates.  Let $b_v$
be the number of colour-block incidences at $v$.  Including the signatures
needed for blocks with no common dart, the explicit vector length is
\[
 O\!\left(\sum_v k_v(h_v+b_v)\right)=O(N^2),
\]
because each of $\sum_v h_v$, $\sum_v k_v$, and $\sum_v b_v$ is $O(N)$.
Lexicographic sorting of these vectors identifies equal region-nodes and
looks up the private-block signatures in $O(N^2\log N)$ time.  The trees
themselves have $\sum_v h_v=O(N)$ incidences.  Ordering their region-node
neighbours by the positions of the distinguished darts in $\omega_v$,
tracing the contours, and inserting the private runs all fit within the
$O(N^3)$ bound.  This proves the time claim in
Theorem~\ref{rep:thm:repair}.

The returned rotation systems specify the full embeddings because both
colours are connected.  If explicit common-component positions are
desired, they can also be produced in polynomial time: for each common
component $X$, search the components of $G_i-V(X)$ and read their faces
of $X$ from attachment corners in the same way.  An isolated $X$ has one
face.  This requires at most $O(N^2)$ additional work and includes every
isolated common vertex.

\subsection{Why the attachment condition cannot be omitted}
\label{rep:counterexample}

\begin{proposition}[A six-vertex obstruction]
\label{rep:prop:incidence-counterexample}
Replacing cycle side bits by arbitrary face labels, while retaining only
common-order consistency and colour-block lobe equalities, gives an
incorrect repair criterion even when the nontrivial common block is
3-connected.
\end{proposition}

\begin{proof}
Let $K$ have vertices $a,b,c,p,q$, the triangle edges $ab,bc,ca$, and the
six edges joining each of $p,q$ to each of $a,b,c$.  Thus $K=K_5-pq$ is
the triangular bipyramid.  Let $H$ be $K$ together with an isolated
vertex $x$, and set
\[
 G_1=H+px,\qquad G_2=H+qx.
\]
Each colour is planar and connected.  The union is obtained by
subdividing the missing edge $pq$ of $K_5$, so it is biconnected.

Prescribe the same embedding of $K$ in both colours.  Their common orders
agree, and the blocks containing $K$ are exactly $K$ itself.  Consequently
there are no colour lobes of $K$, and every lobe equality is vacuous.
Nevertheless, the six faces of $K$ have boundaries
\[
 pab,\quad pbc,\quad pca,\quad qab,\quad qbc,\quad qca.
\]
No face contains both $p$ and $q$.  In a \SEFE{}, the vertex $x$ and the
interiors of $px,qx$ would form a connected set disjoint from $K$, placing
$x$ in a face incident with both vertices, which is impossible.  Equivalently,
for $Z=\{x\}$ the attachment set is $A_K(Z)=\{p,q\}$ and
$L(K,Z)=\varnothing$.
\end{proof}

The obstruction is already present with a fixed embedding of $K$; it is
not an artefact of changing face names while an embedding varies.  In the
valid repair theorem, the new ingredient is precisely the attachment-face
condition.  Once it is imposed, the tree used by the cactus proof survives
as the block--region incidence tree $T_v$.
